\scanpage{230}
\mspage{180}
% L1–3
$b_n\Rightarrow c_n$. On peut supposer que $f$ est une $\mathrm{Hot}$-fibration (et $\unsure{=}$ $\mathrm{Hot}$-parfait si ça nous chante\ldots). \deleted{\unsure{Or}} \deleted{On a la suite spectrale de Leray, pour $f$} \inserted{On sait déjà que \unsure{pour} :}
% L4
\[
H^i(Y,F)\unsure{-|\to} H^i(X,f^*F)
\tag*{($\alpha$)}
\]
% L5–13
tout est OK si $i\in\{0,1\}$. Donc on peut supposer $F$ abélien. Si ça nous avance, on peut supposer, par récurrence (que l'implication $b_{n-1}\Rightarrow c_{n-1}$ est connue, donc) que ($\alpha$) est iso pour $i<n$. Reste à le voir pour $i=n$. \deleted{Notons} On est ramené bien sûr, comme $\pi_0(f)$ est bijectif, au cas où $X$, $Y$ sont 0-connexes, donc il y a \unsure{esst.} une seule fibre $Z$.

% L14–15
La suite exacte d'homotopie nous \deleted{\unsure{dit}} apprend que
% L16–17
\[
\begin{aligned}
\pi_{\unsure{i|1}}(Z,\unsure{z})=0&\qquad\text{pour }i\leq n-1\\
&\qquad(Z\ \text{est }(n-1)\text{-connexe})
\end{aligned}
\]
% L18
et de plus
% L19–21
\[
\begin{aligned}
\pi_n(Z,z)&\longrightarrow\pi_n(X,z)\quad\text{est nul.}\\
\text{(\unsure{car|donc}}\quad\pi_n(X,z)&\to\pi_n(Y,\unsure{y})\quad\text{iso donc mono.)},\\
\text{donc }\pi_{n+1}(Y,\unsure{z|y})&\to\pi_n(Z,z)\quad\text{épi}.
\end{aligned}
\]
% L22
Comme $Z$ est $(n-1)$-connexe, \inserted{(donc \deleted{\unsure{aussi}} 1-connexe)} on a
% L23
\[
H^i(Z,\unsure{G|E|\mathcal G})=0\quad\text{si }0<i<n
\]
% L24
pour $G$ loc. constant (donc constant) sur $Z$.
% archive sheet number: 230
